\documentclass[UTF8,11pt,a4paper,fontset=windows]{ctexart}
\usepackage[margin=20mm,headheight=15pt]{geometry}
\usepackage{booktabs,tabularx,array,enumitem,xcolor,fancyhdr,amsmath}
\usepackage[unicode,colorlinks=true,linkcolor=blue!45!black,urlcolor=blue!45!black]{hyperref}
\definecolor{ink}{HTML}{184E68}
\setlength{\parindent}{0pt}\setlength{\parskip}{6pt}
\setlist{nosep,leftmargin=1.5em,topsep=3pt}
\renewcommand{\arraystretch}{1.28}
\newcolumntype{Y}{>{\raggedright\arraybackslash}X}
\newcolumntype{P}[1]{>{\raggedright\arraybackslash}p{#1}}
\pagestyle{fancy}\fancyhf{}
\fancyhead[L]{\small\color{ink}BrewFS $\times$ AX：云上 Agent 状态管理}
\fancyhead[R]{\small 2026-10-23 / 计划，未执行}
\fancyfoot[C]{\small\thepage}
\ctexset{section={format=\Large\bfseries\color{ink},beforeskip=12pt,afterskip=7pt},
subsection={format=\normalsize\bfseries,beforeskip=8pt,afterskip=4pt}}
\begin{document}
{\huge\bfseries\color{ink}实验要回答的四个问题}\par
{\large 从单个任务速度，走向云上集群的状态管理}\par
\textbf{已知条件：}AX 集成和多 Agent 运行已具备。本文只规划对比实验，没有新增测量结果。完整输入、命令、计时和来源见上级目录的 \texttt{plan.zh.md}、\texttt{baselines.md} 和 \texttt{workloads.json}。

\section*{1\quad 研究目标与明确对照}
管理工作区文件、保存版本和可恢复会话。一个 Agent 应能独立创建、修改、分叉、共享资料、换节点继续和释放空间。进程内存及未完成外部事务不属于本轮恢复范围。

{\small\begin{tabularx}{\linewidth}{@{}P{11mm}Y P{37mm}@{}}
\toprule
问题 & 要问清楚什么 & 主要结果\\
\midrule
RQ1 & 同样资源能管理多少工作区？持续创建是否排队？ & 容量、按时完成率、准备延迟\\
RQ2 & 共享资料、独立修改和保留历史，成本是多少？ & 保存/分发时间、真实空间\\
RQ3 & 存储改进能否让真实 Agent 更早完成任务？ & 固定 16 题正确完成曲线\\
RQ4 & 能否只迁走选中 Agent，其他任务照常运行？ & 可选目标、停顿、邻居影响\\
\bottomrule
\end{tabularx}}

{\small\begin{tabularx}{\linewidth}{@{}P{11mm}P{48mm}Y@{}}
\toprule
方案 & 明确组成 & 不能拿掉的已有能力\\
\midrule
B0 & BrewFS＋当前 AX 集成 & 记录共享基线、私有修改、Packed 实际路径\\
B1 & XFS reflink＋restic & 块共享、增量上传、跨工作区去重\\
B2 & JuiceFS CE＋Redis＋RustFS & 原生 clone、对象共享、原目录重挂\\
B3 & CephFS＋RADOS＋多 MDS & 快照、只读共享、原子卷重挂\\
B4 & EBS gp3＋XFS 原卷换挂载 & 原盘直接移交，不强制复制或先做快照\\
\bottomrule
\end{tabularx}}

\textbf{比较边界。}B0--B3 在受控集群比较；B4 与同云、同区域和同可用区的 B0/B2 另做云实验。不能把本地结果混进云端排名。JuiceFS clone 共享数据块但复制元数据；CephFS 子卷 clone 是异步复制；只读消费者直接共享快照。

\textbf{待证明的优势。}共享多、修改少、状态变化频繁时，是否能用更少元数据、保存和搬运成本服务更多 Agent。不是预设“现有方案都不行”。
\newpage

\section*{2\quad T1：规模和排队（RQ1）}
\subsection*{固定资源与输入}
拟定 4 个客户端，各 8 vCPU / 32 GiB，客户端总文件缓存 32 GiB。控制与存储服务总预算 32 vCPU / 128 GiB，4 台同型号存储节点；总磁盘、网络与后台维护预算公开。允许各方案在预算内调优，正式采集前冻结。

F1：32,768 个 16 KiB 小文件＋128 个 4 MiB 文件，共 1 GiB。每目录不超过 256 项；种子 20261023。输出检查路径、类型、长度和 SHA-256。F1/F2 为规格已定、尚未物化的输入；不生成虚构测量。

\subsection*{容量实验}
\begin{enumerate}
\item 一次提交 N=16、64、256、1024 个独立工作区创建请求。
\item 每工作区读固定 256 个小文件，向大文件区写 256 次不重叠 4 KiB，共 1 MiB；另保存 64 KiB 会话。
\item 120 秒准入截止；成功工作区保持打开至第 180 秒，在 120/180 秒核验后统一释放，不能提前关闭绕过容量限制。
\end{enumerate}
主要报告创建成功数、请求到首次正确读取的 p50/p95、元数据内存和物理空间。所有请求进入分母。各 N 下工作量在系统间相同；N 增加不是固定任务加速实验。

\subsection*{持续压力}
64 个槽位，创建与保存分开测试。按预生成时间表，以 1/4/16/64 次每秒发送；预热 5 分钟、测量 15 分钟、最多收尾 10 分钟。记录计划发送、实际发送、开始、完成时刻。创建完成保存核对后释放槽位；保存请求在同一工作区排队。

主要看按时完成率、正确完成数/秒和队列是否持续增长。系统变慢时不能少发请求，再用变少的后端流量宣称省开销。发生器自身延迟单列。

\subsection*{必须保留的边界}
\begin{itemize}
\item 同一资料版本与 64 个内容不同版本；冷、暖客户端缓存。
\item 旧 checkout 的 owner 限制在 256 附近；实际部署补测 255/256/257。
\item 1024 是拟尝试规模，不是已支持。容量拒绝不能当资格失败删除。
\item 分成多个控制域可另测，但披露域数和总资源，不能冒充单域容量。
\end{itemize}

\textbf{运行合同。}需要跨节点恢复的保存时间必须包括远端数据和会话确认。B1 本地 fsync 不代替备份；B0 平时写远端的成本也计入全任务。当前故障目标只覆盖客户端节点，不外推后端故障保证。
\newpage

\section*{3\quad 状态成本：T2 与 T3（RQ2）}
\subsection*{T2：分叉、独立修改、历史与回收}
从 F1 修改 1 MiB 并保存父状态 P，然后派生 16 个候选。共享原始基线不等于继承父任务已保存的修改。

各候选分别改 1/16/128/512 MiB；固定 4 KiB 写，内容修改只在大文件区。每人向同名路径写不同内容。另删 32 个小文件、改名另 32 个、保存 64 KiB 会话。增加相同字节的分散写，以及 1 MiB 重复覆盖十次。

\textbf{输出：}每个候选最终清单，原始基线和父状态清单；恢复父状态、子状态、丢弃子状态后再读父状态。父/兄弟污染或确认数据丢失直接判失败。

\textbf{主指标：}分叉到首次正确读、保存到远端可恢复的时间；保留所有候选的物理空间、对象数和元数据。共享块只数一次；逻辑目录大小不等于物理占用。

\textbf{历史：}活跃工作区固定 16，另保留 1,000/10,000/100,000 个已关闭但可恢复版本。全检根引用，抽 100 个完整恢复，不能把抽样说成全检。删除一半，观察 60 分钟内释放空间和前台延迟。已删除记录、保留版本、同时挂载数分别统计。

\subsection*{T3：生产资料，其他 Agent 直接共享}
F2：64 个 16 MiB 分片，每条记录 4 KiB，key 为 0--262143。生产者生成数据及定位索引，保存后发布 V1。N=1/16/64 个消费者，每人查询 1024 个 key。

{\small\begin{tabularx}{\linewidth}{@{}P{28mm}Y@{}}
\toprule
场景 & 固定操作\\
\midrule
高共享 & 所有消费者查询同一组 key\\
低共享 & 不同消费者查询互不重叠的 key\\
完整扫描 & 每人顺序读完整 1 GiB，作为局部读反例\\
版本变化 & V2 只改 key 模 100 为 0 的记录；旧任务继续读 V1\\
\bottomrule
\end{tabularx}}

\textbf{对照：}B0/B2/B3 都允许一次构建后直接共享只读版本；B1 每节点只恢复一份。逐消费者复制是补充，不是唯一对手。

\textbf{输出和主指标：}有序 key、内容摘要和版本；从生产开始到全部正确完成、发布到首个结果、所有节点的总传输和物理读取。计入首次构建与发布；半发布不可见。单消费者、全扫描、无重叠时无优势也应报告。
\newpage

\section*{4\quad T4：真实 Agent 是否受益（RQ3）}
\subsection*{任务已经固定，不按结果挑题}
从固定 SWE-bench Verified 数据文件，按每仓库 instance ID 字符串顺序取前四项。16 个独立问题、16 个原始代码版本。选择不看运行时间、修复率或 BrewFS 性能；这是四仓库案例组，不代表全部任务。

{\small\begin{tabularx}{\linewidth}{@{}P{31mm}Y@{}}
\toprule
仓库 & 已冻结任务编号\\
\midrule
Django & 10097、10554、10880、10914\\
Pytest & 10051、10081、10356、5262\\
Scikit-learn & 10297、10844、10908、11310\\
SymPy & 11618、12096、12419、12481\\
\bottomrule
\end{tabularx}}

题目、原始版本、提示词和独立评分材料已保存。镜像 digest 与评分执行校验尚未完成。评分材料和标准答案不进入 Agent 工作区；退出码 0 不能代替正确修复。

\subsection*{同一批 16 题，只改变并发}
并发 C=1/4/16，每 Agent 上限 2 vCPU / 4 GiB，总客户端资源不变。统一 FIFO，批次开始提交全部 16 题。每题启动后 30 分钟；整批固定观察 8 小时，排队计入端到端时间。

拟定同一模型请求名 \texttt{deepseek-v4-flash}、温度 0、最多 100 轮和 100,000 输出 token；运行前登记真实模型路由和执行器支持。不能支持的参数须对所有方案统一修订。供应商限流保留并单列，不归因于文件系统。

\subsection*{主图是正确完成曲线}
从统一提交起，画 0--8 小时正确完成数，另报每题启动后 30 分钟成功。独立评分完成后把结果归到候选提交时刻；评分耗时单列。

所有失败、超时和未评分进入结果表。结束得早但正确数更少不算胜出。工具调用时间与模型等待分别展示，但不能扣掉模型等待后声称消除了随机性。真实任务的共享程度要实际测，不能从仓库名相同推断。

\subsection*{固定预算与重复单位}
校准：B0/B2，C=4，3 批，每批 16 题，共 \textbf{96 次执行}。

正式：B0/B1/B2/B3，C=1/4/16，3 批，每批 16 题，共 \textbf{36 批、576 次执行}。校准不并入正式结果。B4 先承担真实云迁移对照，不全量交叉所有因素。

按题目配对，保留仓库和批次分层。三批先报告原始点与描述性范围，不用候选数或 I/O 数冒充独立样本。同题 16 候选可另测，但不称 16 个独立问题。
\newpage

\section*{5\quad T5：迁移是否足够灵活（RQ4）}
\textbf{核心场景：64 个工作区正在运行，只迁 16 个，其余 48 个继续。}

前 10 轮执行读、改 1 MiB、保存 64 KiB 会话和核对；统一屏障后暂停编号 0、4、8 至 60 的 16 个迁移。迁移并发为 1/4/16；迁移者再做 10 轮。48 个邻居按每 5 秒一轮运行 120 轮，覆盖 10 分钟。等待迁移槽位也保持暂停；不迁移组工作量相同。源节点按每连续 16 个分组，目标为下一节点。

{\small\begin{tabularx}{\linewidth}{@{}P{31mm}Y@{}}
\toprule
路径 & 应给予的合理迁移方式\\
\midrule
B1 & 增量保存后在目标恢复，允许复用已有内容\\
B2 / B3 & 原远端目录或子卷重新挂载，不强制 clone\\
B4 独立盘 & 每 Agent 一盘，只卸载并换挂其原盘\\
B4 节点共享盘 & 整盘搬迁计所有受影响任务；只搬目录另计复制恢复\\
B0 & 按实际部署移交所有权并重开工作区\\
\bottomrule
\end{tabularx}}

恢复截止为暂停后 600 秒，全批截止 900 秒；邻居收尾最多 300 秒，每轮请求期限 120 秒。全部并发组使用同一口径。

云端目标实例预先就绪。主比较同可用区；独立盘拟定 20 GiB gp3，登记 IOPS/吞吐、挂载数、API 配额和费用。B0 也在同云条件下运行。不能用本地 sleep 模拟云控制面。

\subsection*{三个必须一起报告的结果}
\begin{enumerate}
\item \textbf{迁移范围：}CPU/内存足够的候选节点中，多少能不移动其他 Agent 就恢复；支持范围与实际成功比例分开。
\item \textbf{迁移代价：}请求到恢复、实际暂停到首个正确工具结果、16 个全恢复时间，以及相对不迁移增加的总任务时间和重做量。
\item \textbf{邻居影响：}48 个未迁移任务的吞吐和工具 p95，相对配对不迁移组的变化。
\end{enumerate}

EBS 原卷不能跨可用区直接挂载，快照建新卷另列；B0/B2 也必须验证权限、网络和后端可达性。不能把未知节点算作可迁移。

\subsection*{故障和会话}
分别在确认前杀源、确认后杀源、旧源断网再返回写入处注入故障，每点/系统 20 次。确认内容不得丢，旧源不得覆盖新发布状态；租约到期不自动证明旧挂载安全退出。强制 detach 不用于正常迁移组。

附加两题真实会话：pytest-10051、sympy-11618，第 3 个完整工具结果后暂停，恢复剩余预算。保存文件不自动等于保存会话。

\textbf{允许云盘胜出。}同区独立盘可能迁得更快。工作区粒度、分叉共享与状态占用是否更好，需要其他实验独立证明，不能以厂商 API 慢替代机制贡献。
\newpage

\section*{6\quad T6、执行次序与结论边界}
\subsection*{T6：把设计特性与任务结果连起来}
{\small\begin{tabularx}{\linewidth}{@{}P{31mm}Y@{}}
\toprule
内部对照 & 固定项、改变项与需要观察的结果\\
\midrule
共享基线 $\times$ Packed & 固定 16 工作区和 F1；共享/独立基线与 Packed/等价普通索引组成 2$\times$2。看创建、搜索、lower 请求和元数据内存，报告交互。\\
私有修改 & T2 的 1/128 MiB，比较保存变化与完整物化；固定 API、确认点。看保存、PUT、上传和物理空间。\\
对象布局 $(B,q)$ & T3 局部与全读，固定内容/请求/校验，使用实际不同的合法新对象布局。看 S3 读取、物理读取、工具时间是否同时改变。\\
\bottomrule
\end{tabularx}}

目录查询分“先列目录后查名”和“直接随机查名”，含正负查询。布局必须读回实际参数；upper 整对象校验和默认布局未变是反例。缺少同 API 对照就标未实现，不用不同文件系统整套差值冒充单机制效果。

\subsection*{先做什么}
\begin{enumerate}
\item T1 校准：B0/B2 $\times$ 并发 1/4 $\times$ 3 批，每批固定 16，共 \textbf{12 批、192 个工作区生命周期}。
\item 云端 B0/B4 独立盘同区单任务迁移，1 MiB 修改、冷目标，5 次配对。先检验云盘是否已足够。
\item T1 四系统 $\times$ 四规模 $\times$ 五批，80 批；再做固定到达率压力。
\item T2/T3 先固定主参数，T6解释差异；再扩修改量、重用程度、历史与回收。
\item L16 完成镜像和评分校验，96 次校准后执行正式 576 次。
\item T5 扩展到 64 运行、迁 16、留 48，补并发、故障与真实会话。
\end{enumerate}

\subsection*{统计与失败}
确定性每格 5 个独立批次；压力每格 3 时段；真实任务 3 批；故障每点 20 次。它们是预定预算，不是已执行数量。报原始批次点、配对差值和范围；5 批不报批次 p99。大量请求可算每批请求 p95，不能虚增独立样本。失败保留分母，仪器问题标明，原记录不删。

\subsection*{顶会论证应站在哪里}
云盘换挂载是合理的部署基线，能回答“已有远端云盘为什么不够”。是否更快由测量决定；仅证明云 API 等待长不足以证明设计贡献。论文需要同时展示规模、共享和私有状态成本、迁移粒度、真实任务效果，以及内部机制解释。

运行前仍需登记实际版本、镜像、机器/云配额、模型路由、会话恢复接口、基线适配命令和合法布局。AX 多 Agent 运行不再列为待完成。

{\small 官方基线依据、完整操作合同及 FAST 2026 ThinkAhead 的相关实验启发见 \texttt{baselines.md}。该工作研究镜像加载，并非原盘重挂；不能引用成重挂劣势的证据。}
\end{document}

